\documentclass[11pt, oneside]{article} 
\usepackage{geometry}
\geometry{letterpaper} 
\usepackage{graphicx}
	
\usepackage{amssymb}
\usepackage{amsmath}
\usepackage{parskip}
\usepackage{color}
\usepackage{hyperref}

\graphicspath{{/Users/telliott/Dropbox/Github-Math/figures/}}
% \begin{center} \includegraphics [scale=0.4] {gauss3.png} \end{center}

\title{Pentagon}
\date{}

\begin{document}
\maketitle
\Large

%[my-super-duper-separator]

\label{sec:pentagon}

\subsection*{pentagon}

Euclid shows a hexagon circumscribed by a circle through all 6 vertices.
\begin{center} \includegraphics [scale=0.20] {hex1.png} \end{center}

With the same general idea, we can imagine dividing the circle up into five equal arcs whose endpoints form chords which are the sides of a regular pentagon.  (The pentagon is not \emph{easy}.  How to actually do this construction will be shown in the next book).

In the figure below, we have not drawn the actual circle.
\begin{center} \includegraphics [scale=0.3] {pent_central_label.png} \end{center}

We can find the angles in various ways.  Perhaps the simplest is that  the central angle must be 1/5 of the complete circle and $\angle CAD$ (light red triangle) is one-half that, or $36^\circ$.

The base angles $\angle ACD$ and $\angle ADC$ cut two adjacent arcs.  So the base angles are each twice the measure of the apex angle.  $\triangle ACD$ is isosceles with angles in the ratio $1:2:2$.

Since $B$ cuts 3 arcs and $\angle BAD$ cuts 2, it follows that $AD \parallel BC$.

\begin{center} \includegraphics [scale=0.3] {pent0a.png} \end{center}

We can use the $1:2:2$ ratio to draw a triangle similar to $\triangle CAD$ from above.  In the triangle above, the apex angle $\alpha = 36^\circ$ and base angles $\beta$, where $\beta = 2 \alpha$.

Now draw the transversal to make the yellow triangle isosceles, scaled to divide the original side into parts of length $x$ and $1$.

By sum of angles, the obtuse angle of the yellow triangle is $3 \alpha$, its supplementary angle is $2 \alpha$, so both the yellow and the red triangle are isosceles.  The red one is similar to the original with base angles equal to $\beta$.

By similar triangles we have
\[ \frac{x+1}{x} = \frac{x}{1} \]
\[ x^2 = x + 1 \]
The solution to this quadratic equation is usually given the symbol $\phi$, the famous golden ratio.

\begin{center} \includegraphics [scale=0.3] {pent0a.png} \end{center}

Use the quadratic formula to obtain the positive root:
\[ \phi = \frac{1 + \sqrt{5}}{2} \]

\subsection*{problem}

\begin{center} \includegraphics [scale=0.35] {phi_circle.png} \end{center}

Here's another construction in the collection curated by Bogomolny, from John Arioni.  Let the diameter of the circle be $1$, so the rectangle has sides of length $1$ and $2$ and the diameter is $\sqrt{5}$.

The distance from corner to corner is twice the distance from the corner to the center of the circle, or $\sqrt{5}/2$.  Add another $1/2$ to go to the edge of the circle and we have $\phi$.

Let the distance to go the rest of the way be $x$.  We have 
\[ x = \sqrt{5}/2 - 1/2 \]
\[ = \sqrt{5}/2 + 1/2 - 1 \]
\[ = \phi - 1 \]

Finally, recall that 
\[ \phi^2 = \phi + 1 \]
\[ \phi = 1 + \frac{1}{\phi} \]

So $\phi - 1 = 1/\phi$.

$\square$

\end{document}